% Conjectural application.  The proved Hasse-jet results do not depend on it.

\section{Tate-dual obstruction classes}

Put $P=p(p-1)$ and let $i=np+i'$ be in the nondegenerate band. Define
\[
  i^\dagger=p^2+1-k-i
   =(p-2-n)p+(2p+1-k-i').
\]

On every admissible cell considered here, the mod-$p^2$ filtration has the
same weight class as $k+2i$.  We therefore write
\[
 \omega_{p^2}(\theta^if)=k+2i+\delta(\theta^if)P.
\]
This equation defines the integer $\delta(\theta^if)$.  Thus
$\delta=-1$ means that one additional full Hasse drop occurs below
$k+2i$, while $\delta=0$ means that it does not.
Then
\[
 i+i^\dagger=p^2+1-k,
 \qquad
 (k+2i)+(k+2i^\dagger)=2p^2+2.
\]
Moreover, if $c(x)=x^2+(k-1)x$, then
$c(2p+1-k-i')\equiv c(i')\pmod p$, while
\[
 (p-2-n)(p-1-n)\equiv(n+1)(n+2)\pmod p.
\]

The lower Bockstein class $\beta_w(g)$ and its lift-independence are given
by Lemma~\ref{lem:bockstein}. In particular,
\[
  \beta_w(g)=0
  \quad\Longleftrightarrow\quad
  \omega_{p^2}(g)\le w
\]
for an admissible candidate weight $w$.

The perfect Hasse-quotient pairing and theta adjointness are proved in
Propositions~\ref{prop:hasse-perfect-pairing} and
\ref{prop:theta-adjointness}. Put
\[
 w_-=k+2i-P,\qquad w_+=k+2i^\dagger+P.
\]
Proposition~\ref{prop:geometric-bockstein} places the two actual
weight-lowering obstructions in
\[
 Q_{r_-}\quad\hbox{and}\quad Q_{r_+},
 \qquad
 r_-=w_-+P=k+2i,
 \quad
 r_+=w_++P=k+2i^\dagger+2P.
\]
Their weights satisfy
\[
 r_-+r_+=P+2+p(3p-1).
\]
Theorem~\ref{thm:natural-bockstein-pairing} uses a horizontal unit of
weight $p(3p-1)$ to give a perfect pairing
\[
 [\ ,\ ]^{\mathrm{nat}}:
 Q_{r_-}\times Q_{r_+}\longrightarrow\mathbb F_p
\]
as part of a correctly typed family under which theta is skew-adjoint.
The pairing shows that the two obstruction spaces are dual.  It does not
by itself identify the two particular obstruction classes.

\begin{conjecture}[Cross-wedge Bockstein compatibility]
\label{conj:cross-wedge-bockstein}
Assume that $f$ is nonzero modulo $p$, is ordinary at $p$, has
$\omega_p(f)=k$, $4\le k<p-1$, and that both $i$ and $i^\dagger$ lie in
the nondegenerate band. If
\[
 D(i)=0,\qquad D(i^\dagger)=2,
\]
then the two geometric Bockstein classes
\[
 \Gamma_{w_-}(\theta^if)\in Q_{r_-},
 \qquad
 \Gamma_{w_+}(\theta^{i^\dagger}f)\in Q_{r_+}
\]
vanish together:
\[
 \Gamma_{w_-}(\theta^if)=0
 \quad\Longleftrightarrow\quad
 \Gamma_{w_+}(\theta^{i^\dagger}f)=0.
\]
Equivalently,
\[
 \beta_{k+2i-P}(\theta^i f)=0
 \quad\Longleftrightarrow\quad
 \beta_{k+2i^\dagger+P}(\theta^{i^\dagger}f)=0.
\]
\end{conjecture}

The conjecture states only the intrinsic vanishing equivalence needed
below.  A stronger equality through a chosen auxiliary-level pairing would
require an explicit comparison map and a normalization of the horizontal
unit.  Those data have not yet been constructed, so no such equality is
claimed here.

\begin{conjecture}[Exact $D=0$ shifted-conic law]
\label{conj:exact-d0-shifted-conic}
Under the hypotheses of the preceding conjecture,
\[
 \delta(\theta^i f)=-1
 \quad\Longleftrightarrow\quad
 i'^2+(k-1)i'-(n+1)(n+2)\equiv0\pmod p.
\]
Otherwise $\delta(\theta^i f)=0$.
\end{conjecture}

\begin{proposition}[Conditional shifted-conic deduction]
\label{prop:conditional-shifted-conic}
Assume Conjecture~\ref{conj:cross-wedge-bockstein}.  Assume also that the
complementary $D=2$ cell lies in the exact range where
\[
 \Gamma_{w_+}(\theta^{i^\dagger}f)=0
 \quad\Longleftrightarrow\quad
 (i'^\dagger)^2+(k-1)i'^\dagger
   -n^\dagger(n^\dagger+1)\equiv0\pmod p.
\]
Then the vanishing equivalence in
Conjecture~\ref{conj:exact-d0-shifted-conic} holds.
\end{proposition}

\begin{proof}
For the dual index, the proved $D=2$ polynomial is
\[
 (i'^\dagger)^2+(k-1)i'^\dagger
   -n^\dagger(n^\dagger+1).
\]
The two elementary congruences above transform this into
\[
 i'^2+(k-1)i'-(n+1)(n+2).
\]
The additional $D=2$ hypothesis says that the complementary obstruction
vanishes exactly when this polynomial is zero.  The assumed cross-wedge
compatibility transfers this vanishing to $i$.  This gives the claimed
equivalence.  Ahlgren--Raum--Richter \cite[Theorem~C]{ARR2026} motivate
this hypothesis: they prove exact values off the exceptional locus in the
stated ranges and a sharper upper bound on that locus.  Their theorem alone
is not being used here as an exact vanishing equivalence in our coordinates.
\end{proof}
